\section*{Chapitre \ref{ch:b2:deep-time-climate} --- Le climat à travers les temps géologiques}

\begin{solution}{exo:b2:deep-time-climate:1}
(a) $900 + 37100 + 700 + 3 = \qty{38703}{PgC}$, soit environ 65 fois les \qty{591}{PgC} de l'air.
(b) $\num{3.5e6}\times0.012011 = \qty{42000}{PgC}$. (c) $591/42000 = \qty{1.4}{\percent}$.
\end{solution}

\begin{solution}{exo:b2:deep-time-climate:2}
(a) Entrées $11.3 + 3 + 12.3 + 8.1 + 2.0 = \qty{36.7}{Tmol.an^{-1}}$~; sorties $26 + 16 =
\qty{42}{Tmol.an^{-1}}$, égales aux incertitudes près. (b) $(11.3 + 3)/36.7 = 0.39$~: nettement
moins de la moitié~; l'essentiel des entrées est du carbone sédimentaire ancien. (c) $\num{3.5e6}/42 = \qty{8.3e4}{an}$.
\end{solution}

\begin{solution}{exo:b2:deep-time-climate:3}
(a) \ce{Mg2SiO4 + 4 CO2 + 2 H2O -> 2 Mg^{2+} + 4 HCO3^- + SiO2}. (b) Avec \ce{2 Mg^{2+} + 4 HCO3^-
-> 2 MgCO3 + 2 CO2 + 2 H2O}, la réaction nette est \ce{Mg2SiO4 + 2 CO2 -> 2 MgCO3 + SiO2}~: deux moles
de dioxyde de carbone par mole de forstérite. (c) $\qty{1}{PgC} = 10^{15}/12.011 =
\qty{8.33e13}{mol}$ de carbone, ce qui exige \qty{4.16e13}{mol} de forstérite, $4.16\times10^{13}
\times\qty{140.7}{g} = \qty{5.86}{Pg}$, soit près de six milliards de tonnes.
\end{solution}

\begin{solution}{exo:b2:deep-time-climate:4}
Glaciation hirnantienne, de 445 à \qty{443}{Ma}, \omterm{def:b2:deep-time-climate:states}{état climatique glaciaire}~; extension maximale de la glace du Paléozoïque tardif, de 307 à
\qty{299}{Ma}, \omterm{def:b2:deep-time-climate:states}{état climatique glaciaire}~; Turonien, de 94 à \qty{89}{Ma}, \omterm{def:b2:deep-time-climate:states}{état climatique de serre} (le plus chaud du Phanérozoïque)~;
PETM, de 55.9 à \qty{55.7}{Ma}, \omterm{def:b2:deep-time-climate:states}{état climatique de serre}~; EECO, de 53 à \qty{49}{Ma}, \omterm{def:b2:deep-time-climate:states}{état climatique de serre}~; début de la
calotte antarctique, vers \qty{34}{Ma}, début de l'actuel \omterm{def:b2:deep-time-climate:states}{état climatique glaciaire}~; calottes nord,
\qty{2.7}{Ma}, \omterm{def:b2:deep-time-climate:states}{état climatique glaciaire}.
\end{solution}

\begin{solution}{exo:b2:deep-time-climate:5}
(a) $63000/(8.314\times288^2) = \qty{0.0914}{K^{-1}}$, qui s'arrondit à 0.09. (b) $74000/(8.314
\times288^2) = \qty{0.107}{K^{-1}}$. (c) $\ln1.2/0.09 = \qty{2.03}{K}$ et $\ln 1.2/0.107 =
\qty{1.70}{K}$. (d) En régime stationnaire, l'altération doit égaler le dégazage, ce qui fixe $e^{k
\Delta T}$ donc $\Delta T$~; la sensibilité ne détermine que la quantité de dioxyde de carbone qui produit ce
réchauffement.
\end{solution}

\begin{solution}{exo:b2:deep-time-climate:6}
Ici $k_T = 0.0914$. (a) Valeur exacte $\exp[7577.5\times(1/288 - 1/293)] = 1.567$~; approximation
$e^{0.0914\times5} = 1.579$~: \qty{0.8}{\percent} trop élevée. (b) Valeur exacte $e^{2.482} = 11.97$~;
approximation $e^{2.741} = 15.50$~: \qty{30}{\percent} trop élevée. (c) $\ln(\text{rapport}) = (E_{\text
a}/R)[x/T_0^2 - x/(T_0(T_0 + x))] = (E_{\text a}/R)\,x^2/(T_0^2(T_0 + x))$. En posant $\ln1.1$~:
$x^2 = 1.043\,(288 + x)$, donc $x = \qty{17.9}{K}$. Au-delà d'environ \qty{18}{K} de réchauffement, la forme
exponentielle surestime l'altération de plus d'un dixième.
\end{solution}

\begin{solution}{exo:b2:deep-time-climate:7}
(a) $1444/276 = 5.23$, et $3.93\log_2 5.23 = \qty{9.4}{W.m^{-2}}$. (b) Presque exactement les
\qty{9}{W.m^{-2}} du Soleil~: les deux forçages se compensent. (c) $(1444 - 276)/419.5 =
\qty{2.8}{ppm}$ par million d'années. Les \qty{3.4}{ppm} viennent d'une droite ajustée à l'ensemble des
données, qui pèse chaque estimation plutôt que les deux extrémités de la courbe lissée.
\end{solution}

\begin{solution}{exo:b2:deep-time-climate:8}
(a) De $3000/20000 = \qty{0.15}{PgC.an^{-1}}$ à $7000/3000 = \qty{2.3}{PgC.an^{-1}}$, une plage
qui contient les 0.3 à 1.5 cités par le GIEC. (b) $35\times12.011/44.01 = \qty{9.6}{PgC.an^{-1}}$, de 6 à
32 fois les 0.3 à \qty{1.5}{PgC.an^{-1}} du PETM. (c) $3.93\log_2(2000/900) = \qty{4.5}{W.m^{-2}}$,
et $3\log_2 2.22 = \qty{3.5}{\celsius}$~: environ soixante-dix pour cent des \qty{5}{\celsius}. D'autres
rétroactions ou une autre source de réchauffement, comme le méthane, doivent combler le reste.
\end{solution}

\begin{solution}{exo:b2:deep-time-climate:9}
(a) Maintenant $1.2\,W_0e^{k\Delta T} = V = W_0$, donc $\Delta T = -\ln 1.2/0.09 = \qty{-2.03}{K}$.
(b) $278\times2^{-2.03/3} = \qty{174}{ppm}$. (c) Non~: en tout état stationnaire, l'altération égale le
dégazage, qui n'a pas changé. Une altérabilité plus forte est compensée par un climat plus froid.
\end{solution}

\begin{solution}{exo:b2:deep-time-climate:10}
(a) $k = 0.119$~: $\tau = 591\times0.693/(0.08\times0.119\times3\times0.138) = \qty{104}{ka}$.
(b) $S = 8$~: $\tau = \qty{51.5}{ka}$. (c) Les deux rendent le thermostat plus rapide, puisque $\tau \propto
1/(kS)$. $E_{\text a}$ est incertain d'un facteur deux environ, $S$ à ces échelles d'un facteur
comparable (3 contre 5 à \qty{8}{K})~: les deux comptent à peu près autant.
\end{solution}

\begin{solution}{exo:b2:deep-time-climate:11}
(a) $0.08\times5000 = \qty{400}{PgC}$ restent dans l'air~: $C/C_0 = 991/591 = 1.677$, soit
\qty{466}{ppm}~; $\Delta T = 3\log_2 1.677 = \qty{2.24}{K}$. (b) $W - W_0 = 0.138\,(e^{0.09\times2.24}
- 1) = \qty{0.0308}{PgC.an^{-1}}$, soit \qty{22}{\percent} d'altération en plus~; $5000/0.0308 =
\qty{162}{ka}$. (c) $\tau = \qty{137}{ka}$ et $\tau\ln10 = \qty{315}{ka}$. (d) Le réchauffement ne croît que
comme le logarithme du dioxyde de carbone, de sorte qu'un grand excès accroît l'altération moins que ne le suppose la
loi linéaire~; l'élimination démarre plus lentement.
\end{solution}

\begin{solution}{exo:b2:deep-time-climate:12}
(a) Près de $M^*$, $\mathrm d\Delta T/\mathrm dM = Sf/(A_0\ln 2\,(1 + fM^*/A_0)) = Sf/(A_0\ln 2\,
C^*/C_0)$ et $\mathrm dW/\mathrm dM = kW\,\mathrm d\Delta T/\mathrm dM$, avec $W = V$ en ce point~; l'inverse est
$\tau'$. (b) $C^*/C_0 = 2^{2.03/3} = 1.597$~: $\tau' = 137.4\times1.597/1.2 =
\qty{183}{ka}$. (c) Plus lent~: à forte teneur en dioxyde de carbone, chaque pétagramme ajouté réchauffe moins,
car le forçage est logarithmique, de sorte que l'altération y répond moins.
\end{solution}

\begin{solution}{pb:b2:deep-time-climate:1}
\textbf{1.} $1.3 + 1.3 + 4.5 + 4.2 = \qty{11.3}{Tmol.an^{-1}}$, $\qty{0.136}{PgC.an^{-1}}$.
\textbf{2.} $11.5\times0.012011 = \qty{0.138}{PgC.an^{-1}}$.
\textbf{3.} Elles concordent à environ \qty{2}{\percent} près~: l'altération des silicates équilibre le dégazage
volcanique, comme l'exige le \omterm{def:b2:deep-time-climate:carbon-cycle}{cycle long du carbone}.
\textbf{4.} $\num{5.1e18}/0.02897 = \qty{1.76e20}{mol}$ d'air~; $\num{591e15}/12.011 =
\qty{4.92e16}{mol}$ de carbone~; $\num{4.92e16}/\num{1.76e20} = \num{2.80e-4}$, soit \qty{280}{ppm}.
\textbf{5.} $591/0.136 = \qty{4350}{an}$.
\textbf{6.} $1 + 0.08 = 1.08$~: seulement \qty{8}{\percent} de plus.
\textbf{7.} $63000/(8.314\times288^2) = \qty{0.0914}{K^{-1}}$, qui s'arrondit à 0.09.
\textbf{8.} $e^{0.09} = 1.09$~: \qty{9}{\percent} de plus par kelvin.
\textbf{9.} $\Delta T = S\log_2(1 + fM/A_0)$~; sa dérivée en 0 vaut $Sf/(A_0\ln2) =
3\times0.08/(591\times0.693) = \qty{5.86e-4}{K.PgC^{-1}}$.
\textbf{10.} $W_0k\times\num{5.86e-4} = 0.138\times0.09\times\num{5.86e-4} = \qty{7.28e-6}{an^{-1}}$.
\textbf{11.} Son inverse est $A_0\ln2/(fkSW_0)$, le $\tau$ de la proposition, car l'équation
linéarisée s'écrit $\mathrm dM/\mathrm dt = -(\mathrm dW/\mathrm dM)\,M$.
\textbf{12.} Les quatre~: $\tau$ est inversement proportionnel à chacun.
\textbf{13.} $1/\num{7.28e-6} = \qty{1.37e5}{an}$.
\textbf{14.} $137\times0.09/0.119 = \qty{104}{ka}$.
\textbf{15.} $k_T = 0.107$~: $137\times0.09/0.107 = \qty{115}{ka}$.
\textbf{16.} Plus courts d'un facteur 1.5 à 7~: le modèle à une boîte et son $f$ constant rendent le
thermostat trop rapide.
\textbf{17.} $\num{50e6}/\num{1.37e5} = 364$~: le thermostat a eu le temps d'équilibrer le cycle à chaque
étape, de sorte que la baisse a suivi des variations lentes du dégazage ou de l'altérabilité.
\textbf{18.} $16\times0.137 = \qty{2.2}{ppm}$.
\textbf{19.} $0.08\times5000 = \qty{400}{PgC}$; $278\times991/591 = \qty{466}{ppm}$.
\textbf{20.} $3\log_2 1.677 = \qty{2.24}{K}$.
\textbf{21.} $0.138\,(e^{0.202} - 1) = \qty{0.0308}{PgC.an^{-1}}$, \qty{22.3}{\percent}.
\textbf{22.} $137\times\ln10 = \qty{315}{ka}$.
\textbf{23.} L'altération a pu être plus sensible (ruissellement, énergie d'activation plus grande, les
rétroactions lentes d'une sensibilité du système Terre plus grande)~; et le carbone organique continental et marin a pu
repousser et retirer du carbone de l'air plus vite que l'altération des silicates.
\textbf{24.} Dans ce modèle, le \omterm{def:b2:deep-time-climate:thermostat}{thermostat de l'altération des silicates} élimine un excès de carbone avec un
temps caractéristique de décroissance d'environ \textbf{\qty[text-series-to-math=true]{137}{} milliers d'années}.
\end{solution}
